\documentclass{article}
\usepackage{pathfinder-readable}

\title{What Rank-Based Reward Shaping Can Learn from Prediction Guarantees}

\begin{document}
\maketitle

\begin{abstract}
Q surveys guarantees for algorithms that use fallible predictions, while P uses multimodal pairwise judgements
to assign credit during cooperative agent training. The thread asked whether Q's account of prediction error
and cost gives a guarantee for P's learned policy. The peer checks found exact cancellation of P's shaping
signal under complete reward accounting, alongside conditional effects from pooling rewards or masking
inactive agents. They also found that P's stated graph condition does not suffice for its unregularised score
estimate. The verifier ended the thread as a draft because these results locate a useful test, but the reward
routing and finite-training benefit in P remain unverified.
\end{abstract}

\section{Introduction}

Q is a survey of learning-augmented algorithms: methods that use predictions while stating what happens when
those predictions are wrong. It separates a predictor's error from the cost of the decision made with it. It
also asks how guarantees pass from one component to the next and how the cost of obtaining a prediction enters
the account. A model's output needs a task-specific meaning and an error measure before an algorithmic
guarantee can use it \cite{Q}.

P addresses credit assignment in cooperative multi-agent learning, including teams whose active members
change. Its MARS-RA method asks a large multimodal model which of two active agents contributed more, fits
Bradley--Terry scores to the judgements, and turns the scores into reward-shaping potentials. P reports
training outcomes and a trend in which more accurate comparisons and more queries accompany better
performance. Its score-estimation claim concerns a latent preference of the comparator; its Shapley-value
interpretation assumes that this preference matches the agents' true contributions \cite{P}.

The thread first asked whether Q's prediction contract could explain a downstream benefit in P. That would
require a link from comparison error to learned-policy performance, together with the price of making
comparisons and a rule for agents entering or leaving the team. The work narrowed to the point where that link
would have to act: the shaped reward and the critic used during finite training.

\section{What was done}

The peers first checked P's score-estimation premise. They then traced one agent's potential through a whole
episode and through a matched critic. An objection about changing teams forced them to repeat the calculation
with an active-agent reward mask. They next checked what happens if the vector of shaped rewards is pooled
into one team reward. Finally, they used a small policy-gradient example to test whether an accurate ranking
is necessarily a useful training signal, and set these results against Q's requirements for downstream error
and query cost.

P's stated protocol compares every ordered pair in the active set. If \(m_t\) agents are active at a queried
step \(t\), this calls the comparator \(m_t(m_t-1)\) times. P reports that, in its experimental setup,
training for 50 million environment steps took 16 hours for MAPPO and 30 hours for MARS-RA. These figures
motivate a charged-query comparison; they do not give a return-per-query guarantee \cite{P}. Q's cost
convention requires a declared conversion from calls to the objective being measured \cite{Q}.

\section{Results}

\subsection{What the potential does exactly}

Let \(i\) denote one member of a fixed agent population, \(t\) a time step, \(T\) the end of an episode, and
\(\gamma\) the discount factor. Write \(\psi_{i,t}\) for that agent's potential at step \(t\), using the same
realised value on the two transitions that meet there. P's shaping increment is
\(F_{i,t}=\gamma\psi_{i,t+1}-\psi_{i,t}\). If every transition is credited and the terminal potential
\(\psi_{i,T}\) is zero, the peers checked the pathwise identity
\[
\sum_{t=0}^{T-1}\gamma^t F_{i,t}
=-\psi_{i,0}+\gamma^T\psi_{i,T}
=-\psi_{i,0}.
\]
The intermediate comparison scores cancel. When the initial potential is fixed before the first action,
shaping shifts an exact episode return by a fixed amount and leaves its policy ordering unchanged. This is an
application of existing potential-shaping results, including work on changing potentials and their relation to
value initialisation \cite{dynamicArchive,dynamicConference,wiewiora}.

The same cancellation reaches an exactly matched critic. Let \(\rho\) be P's shaping weight, \(\widehat
V_i^{\mathrm{env}}\) an unshaped value estimate, and \(\widetilde V_i=\widehat V_i^{\mathrm{env}}-\rho\psi_i\)
its matched shaped version. The peers checked that their one-step temporal-difference residuals are equal. If
the shaped critic has an additional error \(e_i\), the shaped residual differs by \(\gamma
e_{i,t+1}-e_{i,t}\). Thus the comparison signal has no effect through an exactly matched residual, but critic
mismatch gives a specific channel to measure during approximate training. This calculation does not show how
P's implementation represents its critic.

\subsection{What changes with a changing team}

Complete accounting matters. Let \(a_{i,t}\) equal one when agent \(i\) is active at \(t\) and zero otherwise.
If the shaping increment is paid only when the agent is active, the peers obtained, with a zero terminal
potential,
\[
\sum_{t=0}^{T-1}\gamma^t a_{i,t}F_{i,t}
=-a_{i,0}\psi_{i,0}
+\sum_{t=1}^{T-1}\gamma^t(a_{i,t-1}-a_{i,t})\psi_{i,t}.
\]
On re-entry, the final sum can contain an uncompensated negative term. Setting the potential to zero while an
agent is absent does not remove it. P gives a potential over the current active set but does not specify
enough about inactive-agent rewards and critic targets to determine whether its training uses this mask
\cite{P}.

Pooling gives another conditional result. P forms potentials by taking a softmax of the active agents' scores,
so their potentials sum to one whenever the active set is nonempty. If all agent coordinates, including zero
coordinates for absent agents, are summed before a common team update, the pooled shaping increment depends
only on this sum and carries no individual ranking information. Pooling only agents active at the start of a
transition can retain a rank effect at entry, where the boundary accounting above also matters. Separate
per-agent rewards can affect approximate updates. The peers established these alternatives from P's equations,
but did not establish which route its implementation takes.

\subsection{Why a better ranking is not enough}

The peers also checked a one-step diagnostic. One agent chooses a binary action \(a\), with each action
equally likely; a second agent is inert, and the team return is \(G=a\). For a logit policy, the score is
\(z=a-1/2\). With a state-only baseline \(b\), the gradient sample \(g_b=z(G-b)\) has variance
\((b-1/2)^2/4\). The neutral baseline \(b=1/2\) therefore has zero variance, while the softmax of the
rank-correct Shapley scores \((1/2,0)\) gives the first agent \(b>1/2\) and positive variance. This tests
whether ordinal accuracy alone predicts gradient quality; it is neither a MAPPO result nor a contradiction of
P's reported experiments. Prior work already shows that potential shaping can raise policy-gradient variance
\cite{gupta}. Work on preference queries also identifies a gap between query choice and policy improvement
\cite{preferenceQueries}.

\section{Objections and their status}

The first objection was to P's finite-sample Bradley--Terry claim. With two agents at a true tie and two
independent comparisons, the comparison graph is connected. Both outcomes favour the same agent with
probability \(1/2\). On those samples an unregularised likelihood has no finite maximiser: the favoured score
can keep moving away. This exceeds the failure probability \(1/4\) allowed by P's stated \(1-n^{-2}\) claim
when the number of agents \(n\) is two. The peers agreed on this counterexample. A constrained or regularised
estimator and suitable sampling assumptions would be needed for a repaired result; the calculation does not
dispute P's empirical outcomes \cite{P}.

The second objection was to a simpler cancellation argument. Zero potentials at absence and termination are
insufficient when shaping rewards are masked for inactive agents. The active-mask identity above settled the
mathematical point, while the implementation question remained open. The peers also noted that repeated model
calls need not be independent, that P's comparison accuracy does not by itself control critic error, and that
cancellation and variance effects have prior literature. No improved-return theorem or novelty claim follows
from these checks.

\section{Why the thread stopped}

The verifier marked the thread DRAFT because it had reached a conditional result with a clear limit. Complete
shaping cancels from exact returns; full pooling removes the ranking from a common reward; active-only masking
can instead leave a re-entry term. These facts explain why Q's prediction-error framework cannot yet be
transferred to P's learned-policy return. The smallest step to restart the work is to trace one entry or
re-entry transition through P's reward tensor, inactive-agent mask, and critic target. A matched-critic and
neutral-potential control could then test whether any remaining effect improves held-out unshaped return under
stated step, time, and query budgets.

\bibliographystyle{plainurl}
\bibliography{references}

\end{document}
